\documentclass{notaaula}

        \titulonota{Luz, reflexão e refração}
        \creditos{Turma 2B · Física · Setembro/2026 · Prof. Josué Cavalcante}

        \begin{document}
        \cabecalho

        \begin{sobre}
        Introdução à óptica geométrica: fontes de luz, velocidade da luz, propagação, reflexão e refração. Habilidades em foco: EM13CNT102, EM13CNT104, EM13CNT303 e EM13CNT306.
        \end{sobre}

        \section{Luz e óptica geométrica}

\begin{copiar}{Modelo de raios}
A óptica geométrica representa a propagação da luz por raios. Em meio homogêneo e transparente,
eles se propagam aproximadamente em linha reta. No vácuo,
\[
  c\approx3\times10^8\un{m/s}.
\]
Uma fonte \dest{primária} emite luz; uma fonte \dest{secundária} apenas reflete luz recebida.
Feixes podem ser paralelos, convergentes ou divergentes.
\diaadia{Sombras, eclipses e câmeras escuras são explicados pela propagação retilínea.}
\end{copiar}

\section{Reflexão}

\begin{copiar}{Leis da reflexão}
O raio incidente, a normal e o raio refletido pertencem ao mesmo plano. Os ângulos são medidos
em relação à normal e obedecem a
\[
  \theta_i=\theta_r.
\]
Reflexão regular ocorre em superfície lisa e preserva a organização dos raios; reflexão difusa
espalha os raios e permite enxergar objetos de muitos pontos.
\atencao{Não meça o ângulo a partir da superfície. A referência é a normal, perpendicular a ela.}
\end{copiar}

\begin{exemplo}
Um raio forma $25^\circ$ com a normal de um espelho plano. O ângulo de reflexão vale
\resultado{$25^\circ$}; o ângulo entre o raio e o espelho vale $65^\circ$.
\end{exemplo}

\begin{copiar}{Espelho plano}
A imagem em espelho plano é virtual, direita, do mesmo tamanho e fica à mesma distância atrás
do espelho que o objeto está à frente. Há inversão de orientação frente--trás, percebida como
inversão lateral na leitura de palavras.
\end{copiar}

\section{Refração}

\begin{copiar}{Índice e lei de Snell}
Ao mudar de meio, a frequência da luz permanece; a rapidez e o comprimento de onda podem mudar.
\[
  n=\frac{c}{v},
  \qquad
  n_1\sen\theta_1=n_2\sen\theta_2.
\]
Ao entrar em meio de maior índice, o raio se aproxima da normal. Ao entrar em meio de menor índice,
afasta-se da normal.
\simbolos{$n$ não tem unidade; $v$ é a rapidez da luz no meio.}
\end{copiar}

\begin{exemplo}[Rapidez no vidro]
Para um vidro de índice $n=\dec{1,5}$,
\[
  v=\frac{c}{n}=\frac{3\times10^8}{\dec{1,5}}
  =\resultado{2\times10^8\un{m/s}}.
\]
\end{exemplo}

\begin{copiar}{Reflexão total}
A reflexão interna total só pode ocorrer quando a luz tenta passar de maior para menor índice.
Acima do ângulo limite $L$, não há raio refratado:
\[
  \sen L=\frac{n_{\mathrm{menor}}}{n_{\mathrm{maior}}}.
\]
Fibras ópticas conduzem luz por sucessivas reflexões internas totais.
\end{copiar}

\begin{dica}
A frequência é definida pela fonte e não muda na passagem de um meio para outro. Como $v=\lambda f$,
se $v$ muda e $f$ não, então $\lambda$ muda.
\end{dica}

\begin{exercicios}
\nivel{1}{Conceitos}
\begin{questoes}
\item Diferencie fonte primária de fonte secundária.
\item Qual é a rapidez aproximada da luz no vácuo?
\item Um raio incide a $40^\circ$ da normal. Qual é o ângulo de reflexão?
\item Dê duas características da imagem em espelho plano.
\item O que permanece constante quando a luz refrata: frequência ou rapidez?
\nivel{2}{Aplicação}
\item Em um meio com $n=\dec{1,2}$, calcule a rapidez da luz.
\item Um objeto está a $\dec{0,80}\un{m}$ de um espelho plano. Qual é a distância objeto--imagem?
\item A luz passa do ar para o vidro. Ela se aproxima ou se afasta da normal?
\item Calcule $\sen L$ para passagem de vidro ($n=\dec{1,5}$) para ar ($n=1$).
\nivel{3}{Síntese}
\item Explique por que uma piscina parece menos profunda do que realmente é.
\item Um raio no ar incide com $30^\circ$ em vidro de $n=\dec{1,5}$. Use $\sen30^\circ=\dec{0,5}$ e determine $\sen\theta_2$.
\item Relacione reflexão interna total e transmissão de dados em fibras ópticas.
\end{questoes}
\gabarito{3) $40^\circ$; 5) frequência; 6) $\dec{2,5}\times10^8\un{m/s}$;
7) $\dec{1,6}\un{m}$; 8) aproxima-se; 9) $2/3$; 11) $1/3$.}
\end{exercicios}


\end{document}
